%%% Zahlen und Rechenarten
\begin{lesen}
\lies{$12 + 30 = 42$}{12|じゅうに と 30|さんじゅう の 和|わ は 42|よんじゅうに で ある 。}{Die Summe von zwölf und dreißig ist zweiundvierzig.}
\lies{$3 \mid 12$}{3|さん は 12|じゅうに の 約数|やくすう で ある 。 12|じゅうに は 3|さん の 倍数|ばいすう で ある 。 12|じゅうに は 3|さん で 割り切れる|わりきれる 。}{Drei ist ein Teiler von zwölf; zwölf ist ein Vielfaches von drei; zwölf ist durch drei teilbar.}
\lies{$\operatorname{ggT}(12;\,18) = 6$}{12|じゅうに と 18|じゅうはち の 最大|さいだい 公約数|こうやくすう は 6|ろく で ある 。}{Der größte gemeinsame Teiler von zwölf und achtzehn ist sechs.}
\lies{$17$}{17|じゅうなな は 素数|そすう で ある 。 約数|やくすう は 1|いち と その 数|かず 自身|じしん の 2つ|ふたつ だけ で ある 。}{Siebzehn ist eine Primzahl: Sie hat genau zwei Teiler, die Eins und sich selbst.}
\lies{$3{,}14159 \approx 3{,}14$}{小数|しょうすう 第|だい 3|さん 位|い を 四捨五入|ししゃごにゅう して 、 小数|しょうすう 第|だい 2|に 位|い まで 求めなさい|もとめなさい 。}{Runde die Zahl auf zwei Nachkommastellen.}
\lies{$|-7| = 7$}{マイナス 7|なな の 絶対値|ぜったいち は 7|なな で ある 。}{Der Betrag von minus sieben ist sieben.}
\lies{$25\,\%$ von $80$ sind $20$}{80|はちじゅう の 25|にじゅうご パーセント は 20|にじゅう で ある 。}{Fünfundzwanzig Prozent von achtzig sind zwanzig.}
\end{lesen}

%%% Algebra, Gleichungen und Mengen
\begin{lesen}
\lies{$3(x + 2) = 3x + 6$}{かっこ を はずして 展開|てんかい しなさい 。}{Multipliziere die Klammer aus.}
\lies{$6x + 9 = 3(2x + 3)$}{共通|きょうつう 因数|いんすう 3|さん で くくりなさい 。}{Klammere aus: Der gemeinsame Faktor ist drei.}
\lies{$(x + 3)^2 = x^2 + 6x + 9$}{かっこ $x$ プラス 3|さん かっことじ の 2乗|にじょう イコール $x$ 2乗|にじょう プラス 6|ろく $x$ プラス 9|きゅう 。}{$x$ plus drei in Klammern zum Quadrat ist gleich $x$ Quadrat plus sechs $x$ plus neun.}
\lies{$2x + 3 = 11$ \newline $x = 4$}{方程式|ほうていしき を $x$ に ついて 解きなさい|ときなさい 。 解|かい は $x$ イコール 4|よん で ある 。}{Löse die Gleichung nach $x$ auf. Die Lösung ist $x$ gleich vier.}
\lies{$x + y = 5$ \newline $x - y = 1$}{連立|れんりつ 方程式|ほうていしき を 解きなさい|ときなさい 。 解|かい は $x$ イコール 3|さん 、 $y$ イコール 2|に で ある 。}{Löse das Gleichungssystem: $x$ gleich drei und $y$ gleich zwei.}
\lies{$A = \{1;\,2;\,3\},\ B = \{3;\,4\}$ \newline $A \cap B = \{3\}$}{$A$ と $B$ の 共通|きょうつう 部分|ぶぶん の 要素|ようそ は 3|さん だけ で ある 。}{Die Schnittmenge von $A$ und $B$ enthält nur das Element drei.}
\end{lesen}

%%% Funktionen
\begin{lesen}
\lies{$f(x) = 2x - 4$}{$f$ は 傾き|かたむき 2|に 、 $y$ 切片|せっぺん マイナス 4|よん の 一次|いちじ 関数|かんすう で ある 。 零点|れいてん は $x$ イコール 2|に で ある 。}{$f$ ist eine lineare Funktion mit der Steigung zwei und dem $y$-Achsenabschnitt minus vier. Ihre Nullstelle ist $x$ gleich zwei.}
\lies{$f(x) = (x - 1)^2 - 4$}{グラフ は 頂点|ちょうてん が $S(1,\,-4)$ の 放物線|ほうぶつせん で ある 。 零点|れいてん は マイナス 1|いち と 3|さん で ある 。}{Der Graph ist eine Parabel mit dem Scheitelpunkt $S(1 \mid -4)$. Die Nullstellen sind minus eins und drei.}
\lies{$f(x) = x^2$}{グラフ は $y$ 軸|じく に 関して|かんして 対称|たいしょう で ある 。 関数|かんすう は $-\infty < x \le 0$ で 減少|げんしょう し 、 $0 \le x < \infty$ で 増加|ぞうか する 。}{Der Graph ist achsensymmetrisch zur $y$-Achse. Die Funktion ist fallend für $-\infty < x \le 0$ und steigend für $0 \le x < \infty$.}
\lies{$f(x) = 2^x$}{この 指数|しすう 関数|かんすう の グラフ は 点|てん $(0,\,1)$ を 通り|とおり 、 $x$ 軸|じく が 漸近線|ぜんきんせん で ある 。}{Der Graph dieser Exponentialfunktion verläuft durch den Punkt $(0 \mid 1)$; die $x$-Achse ist Asymptote.}
\lies{$D_f = \R$ \newline $W_f = \{y \in \R \mid y \ge -4\}$}{定義域|ていぎいき は 実数|じっすう 全体|ぜんたい 、 値域|ちいき は マイナス 4|よん 以上|いじょう の 実数|じっすう 全体|ぜんたい で ある 。}{Der Definitionsbereich ist die Menge aller reellen Zahlen; der Wertebereich ist die Menge aller reellen Zahlen, die größer oder gleich minus vier sind.}
\end{lesen}

%%% Geometrie: Grundbegriffe, Winkel, Abbildungen
\begin{lesen}
\lies{$\overline{AB} = 5\,\mathrm{cm}$}{線分|せんぶん $AB$ の 長さ|ながさ は 5|ご センチメートル で ある 。}{Die Strecke $A$ $B$ ist fünf Zentimeter lang.}
\lies{$g \parallel h$}{点|てん $P$ を 通り|とおり 、 直線|ちょくせん $g$ に 平行|へいこう な 直線|ちょくせん $h$ を ひきなさい 。}{Zeichne durch den Punkt $P$ die Parallele $h$ zu $g$.}
\lies{$\alpha = 35\degree$}{アルファ は 35|さんじゅうご 度|ど の 鋭角|えいかく で ある 。}{Alpha ist ein spitzer Winkel von fünfunddreißig Grad.}
\lies{$\beta = 180\degree - 70\degree = 110\degree$}{隣接|りんせつ 補角|ほかく の 和|わ は 180|ひゃくはちじゅう 度|ど で ある 。}{Nebenwinkel ergänzen sich zu hundertachtzig Grad.}
\lies{\skSpiegelung}{三角形|さんかくけい を 、 直線|ちょくせん $g$ に 関して|かんして 対称|たいしょう 移動|いどう しなさい 。}{Spiegele das Dreieck an der Geraden $g$.}
\lies{\skDrehung}{三角形|さんかくけい を 、 点|てん $Z$ を 中心|ちゅうしん に 90|きゅうじゅう 度|ど 回転|かいてん 移動|いどう しなさい 。}{Drehe das Dreieck um neunzig Grad um den Punkt $Z$.}
\lies{\skMittelsenkrechte}{線分|せんぶん $AB$ の 垂直|すいちょく 二等分線|にとうぶんせん を 作図|さくず しなさい 。}{Konstruiere die Mittelsenkrechte der Strecke $A$ $B$.}
\end{lesen}

%%% Ebene Figuren
\begin{lesen}
\lies{$a = 6\,\mathrm{cm},\ b = 4\,\mathrm{cm}$ \newline $A = 24\,\mathrm{cm}^2$}{長方形|ちょうほうけい の 面積|めんせき は 、 6|ろく センチメートル 掛ける|かける 4|よん センチメートル で 、 24|にじゅうよん 平方|へいほう センチメートル で ある 。}{Der Flächeninhalt des Rechtecks ist sechs Zentimeter mal vier Zentimeter, also vierundzwanzig Quadratzentimeter.}
\lies{$3^2 + 4^2 = 5^2$}{3辺|さんぺん の 長さ|ながさ が 3|さん 、 4|よん 、 5|ご の 三角形|さんかくけい は 、 直角|ちょっかく 三角形|さんかくけい で ある 。}{Ein Dreieck mit den Seitenlängen drei, vier und fünf ist rechtwinklig.}
\lies{\skGleichschenklig}{二等辺|にとうへん 三角形|さんかくけい で は 、 2辺|にへん が 等しく|ひとしく 、 2つ|ふたつ の 底角|ていかく も 等しい|ひとしい 。}{Ein gleichschenkliges Dreieck hat zwei gleich lange Seiten und zwei gleich große Basiswinkel.}
\lies{\skDiagonale}{長方形|ちょうほうけい の 対角線|たいかくせん は 長さ|ながさ が 等しく|ひとしく 、 それぞれ の 中点|ちゅうてん で 交わる|まじわる 。}{Die Diagonalen eines Rechtecks sind gleich lang und halbieren einander.}
\lies{$d = 10\,\mathrm{cm}$ \newline $u = \pi \cdot d \approx 31{,}4\,\mathrm{cm}$}{直径|ちょっけい 10|じゅう センチメートル の 円|えん の 円周|えんしゅう は 、 約|やく 31.4|さんじゅういってん_よん センチメートル で ある 。}{Der Umfang eines Kreises mit zehn Zentimetern Durchmesser beträgt ungefähr einunddreißig Komma vier Zentimeter.}
\end{lesen}

%%% Körper
\begin{lesen}
\lies{$a = 3\,\mathrm{cm}$ \newline $V = a^3 = 27\,\mathrm{cm}^3$}{1辺|いっぺん が 3|さん センチメートル の 立方体|りっぽうたい の 体積|たいせき は 、 27|にじゅうなな 立方|りっぽう センチメートル で ある 。}{Das Volumen eines Würfels mit drei Zentimetern Kantenlänge beträgt siebenundzwanzig Kubikzentimeter.}
\lies{$O = 6a^2 = 54\,\mathrm{cm}^2$}{表面積|ひょうめんせき は 54|ごじゅうよん 平方|へいほう センチメートル で ある 。}{Sein Oberflächeninhalt beträgt vierundfünfzig Quadratzentimeter.}
\lies{\skNetz}{立方体|りっぽうたい の 展開図|てんかいず は 、 6つ|むっつ の 正方形|せいほうけい から なる 。}{Das Netz eines Würfels besteht aus sechs Quadraten.}
\lies{\skPyramide}{正四角錐|せいしかくすい に は 、 頂点|ちょうてん が 5つ|いつつ 、 辺|へん が 8本|はっぽん 、 面|めん が 5つ|いつつ ある 。}{Eine quadratische Pyramide hat fünf Ecken, acht Kanten und fünf Flächen.}
\lies{$r = 2\,\mathrm{cm},\ h = 5\,\mathrm{cm}$ \newline $V = \pi r^2 h \approx 62{,}8\,\mathrm{cm}^3$}{円柱|えんちゅう の 体積|たいせき は 20|にじゅう パイ で 、 約|やく 62.8|ろくじゅうに_てん_はち 立方|りっぽう センチメートル で ある 。}{Das Volumen des Zylinders ist zwanzig pi, ungefähr zweiundsechzig Komma acht Kubikzentimeter.}
\end{lesen}

%%% Analysis
\begin{lesen}
\lies{$f(x) = x^3 - 3x$ \newline $f'(x) = 3x^2 - 3$}{$f(x)$ を 微分|びぶん しなさい 。 導関数|どうかんすう は 3|さん $x$ 2乗|にじょう マイナス 3|さん で ある 。}{Leite $f$ ab: Die Ableitung ist drei $x$ Quadrat minus drei.}
\lies{$f'(x) = 0 \ \Rightarrow\ x = \pm 1$}{導関数|どうかんすう は 、 $x$ イコール プラス マイナス 1|いち で 0|ゼロ に なる 。}{Die Ableitung ist null für $x$ gleich plus oder minus eins.}
\lies{$H(-1 \mid 2),\ T(1 \mid -2)$ \newline $W(0 \mid 0)$}{グラフ は 点|てん $(-1,\,2)$ で 極大|きょくだい 、 点|てん $(1,\,-2)$ で 極小|きょくしょう と なり 、 原点|げんてん が 変曲点|へんきょくてん で ある 。}{Der Graph hat den Hochpunkt $H$, den Tiefpunkt $T$ und den Wendepunkt $W$ im Ursprung.}
\lies{$F(x) = x^2$ \newline $f(x) = 2x$}{$F$ は $f$ の 原始|げんし 関数|かんすう で ある 。}{Groß $F$ ist eine Stammfunktion von $f$.}
\end{lesen}

%%% Vektoren und analytische Geometrie
\begin{lesen}
\lies{$\vec{a} = \vektor{3}{4}{0}$ \newline $|\vec{a}\,| = \sqrt{9 + 16 + 0} = 5$}{ベクトル $a$ の 大きさ|おおきさ は 5|ご で ある 。}{Der Betrag des Vektors $a$ ist fünf.}
\lies{$A(1 \mid 2 \mid 3),\ B(4 \mid 6 \mid 3)$ \newline $\overrightarrow{AB} = \vektor{3}{4}{0}$}{点|てん $A$ から 点|てん $B$ へ の ベクトル の 成分|せいぶん は 、 3|さん 、 4|よん 、 0|ゼロ で ある 。}{Der Vektor von $A$ nach $B$ ist drei, vier, null.}
\lies{$g\colon\ \vec{x} = \vektor{1}{2}{3} + t \cdot \vektor{3}{4}{0}$}{直線|ちょくせん $g$ ： 位置|いち ベクトル プラス $t$ 掛ける|かける 方向|ほうこう ベクトル}{die Gerade $g$: Stützvektor plus $t$ mal Richtungsvektor}
\end{lesen}

%%% Stochastik
\begin{lesen}
\lies{$P(\text{„6“}) = \dfrac16$}{正しく|ただしく 作られた|つくられた さいころ を 投げて|なげて 6|ろく の 目|め が 出る|でる 確率|かくりつ は 、 6|ろく 分|ぶん の 1|いち で ある 。}{Die Wahrscheinlichkeit, mit einem fairen Würfel eine Sechs zu würfeln, beträgt ein Sechstel.}
\lies{$P(\text{„zweimal Kopf“}) = \dfrac12 \cdot \dfrac12 = \dfrac14$}{硬貨|こうか を 2回|にかい 投げる|なげる 。 2回|にかい とも 表|おもて が 出る|でる 確率|かくりつ は 、 4|よん 分|ぶん の 1|いち で ある 。}{Eine faire Münze wird zweimal geworfen. Die Wahrscheinlichkeit für zweimal Kopf beträgt ein Viertel.}
\lies{$P(\overline{A}) = 1 - P(A)$}{余事象|よじしょう の 確率|かくりつ は 、 1|いち 引く|ひく $A$ の 確率|かくりつ に 等しい|ひとしい 。}{Die Wahrscheinlichkeit des Gegenereignisses ist eins minus die Wahrscheinlichkeit von $A$.}
\lies{$2;\ 4;\ 6;\ 8$ \newline $\bar{x} = 5$}{4つ|よっつ の 値|あたい の 平均値|へいきんち は 5|ご で ある 。}{Der Mittelwert der vier Werte ist fünf.}
\lies{$\E(X) = 3{,}5$}{さいころ の 出る|でる 目|め の 期待値|きたいち は 3.5|さん_てん_ご で ある 。}{Der Erwartungswert der Augenzahl eines fairen Würfels ist drei Komma fünf.}
\end{lesen}

%%% Logik und Beweis
\begin{lesen}
\lies{$x > 2 \ \Rightarrow\ x^2 > 4$}{$x$ が 2|に より 大きい|おおきい ならば 、 $x$ 2乗|にじょう は 4|よん より 大きい|おおきい 。}{Wenn $x$ größer als zwei ist, dann ist $x$ Quadrat größer als vier.}
\lies{$n$ gerade $\Leftrightarrow$ $n^2$ gerade}{$n$ が 偶数|ぐうすう で ある こと と 、 $n$ 2乗|にじょう が 偶数|ぐうすう で ある こと は 同値|どうち で ある 。}{$n$ ist genau dann gerade, wenn $n$ Quadrat gerade ist.}
\lies{$\forall\, x \in \R\colon\ x^2 \ge 0$}{すべて の 実数|じっすう $x$ に ついて 、 $x$ 2乗|にじょう は 0|ゼロ 以上|いじょう で ある 。}{Für alle reellen Zahlen $x$ gilt: $x$ Quadrat ist größer oder gleich null.}
\lies{$\exists\, x \in \R\colon\ x^2 = 2$}{$x$ 2乗|にじょう イコール 2|に と なる 実数|じっすう $x$ が 存在|そんざい する 。}{Es gibt eine reelle Zahl $x$, für die $x$ Quadrat gleich zwei ist.}
\lies{Annahme}{ルート 2|に が 有理数|ゆうりすう で ある と 仮定|かてい する 。 \ldots\ これ は 矛盾|むじゅん で ある 。 （ 背理法|はいりほう ）}{Angenommen, die Wurzel aus zwei ist rational. \ldots\ Das ist ein Widerspruch. (Beweis durch Widerspruch)}
\end{lesen}

%%% Zahlen, Zeichen und Schreibweisen
\begin{lesen}
\lies{$7{,}25$}{7.25|なな_てん_に_ご}{sieben Komma zwei fünf}
\lies{$0{,}05$}{0.05|れい_てん_れい_ご}{null Komma null fünf}
\lies{$1\,250\,000$ \newline \footnotesize jap. 1,250,000}{125万|ひゃくにじゅうごまん}{eine Million zweihundertfünfzigtausend}
\end{lesen}

%%% Relationszeichen
\begin{lesen}
\lies{$5 < 8$}{5|ご は 8|はち より 小さい|ちいさい}{fünf ist kleiner als acht}
\lies{$x \ge 0$}{$x$ は 0|ゼロ 以上|いじょう ： $x$ は 負|ふ で ない}{$x$ ist größer oder gleich null: $x$ ist nicht negativ}
\lies{$7 \neq 9$}{7|なな は 9|きゅう に 等しくない|ひとしくない}{sieben ist ungleich neun}
\end{lesen}

%%% Rechnen mit Brüchen
\begin{lesen}
\lies{$\dfrac56 - \dfrac14 = \dfrac{10}{12} - \dfrac{3}{12} = \dfrac{7}{12}$}{6|ろく 分|ぶん の 5|ご 引く|ひく 4|よん 分|ぶん の 1|いち は 、 12|じゅうに 分|ぶん の 10|じゅう 引く|ひく 12|じゅうに 分|ぶん の 3|さん で 、 12|じゅうに 分|ぶん の 7|なな}{fünf Sechstel minus ein Viertel ist gleich zehn Zwölftel minus drei Zwölftel, also sieben Zwölftel}
\lies{$\dfrac45 \cdot 10 = 8$}{10|じゅう の 5|ご 分|ぶん の 4|よん は 8|はち}{vier Fünftel von zehn sind acht}
\lies{$2\tfrac12 = \dfrac52$}{2|に と 2|に 分|ぶん の 1|いち は 、 2|に 分|ぶん の 5|ご}{zweieinhalb ist gleich fünf Halbe}
\lies{$\dfrac38 = 0{,}375 = 37{,}5\,\%$}{8|はち 分|ぶん の 3|さん は 0.375|れい_てん_さん_なな_ご 、 すなわち 37.5|さんじゅうなな_てん_ご パーセント}{drei Achtel ist gleich null Komma drei sieben fünf oder siebenunddreißig Komma fünf Prozent}
\end{lesen}

%%% Potenzen, Wurzeln und Logarithmen
\begin{lesen}
\lies{$2^{10} = 1024$}{2|に の 10乗|じゅうじょう は 1024|せんにじゅうよん}{zwei hoch zehn ist eintausendvierundzwanzig}
\lies{$a^m \cdot a^n = a^{m+n}$}{$a$ の $m$ 乗|じょう 掛ける|かける $a$ の $n$ 乗|じょう は 、 $a$ の $m$ プラス $n$ 乗|じょう}{$a$ hoch $m$ mal $a$ hoch $n$ ist gleich $a$ hoch $m$ plus $n$}
\lies{$9^{\frac12} = \sqrt{9} = 3$}{9|きゅう の 2|に 分|ぶん の 1|いち 乗|じょう は ルート 9|きゅう で 、 3|さん}{neun hoch ein Halb ist gleich der Wurzel aus neun, also drei}
\lies{$\sqrt[3]{27} = 3$}{27|にじゅうなな の 3乗根|さんじょうこん は 3|さん}{die dritte Wurzel aus siebenundzwanzig ist drei}
\end{lesen}

%%% Zahlbereiche
\begin{lesen}
\lies{$-3 \in \Z,\ -3 \notin \N$}{マイナス 3|さん は 整数|せいすう で ある が 、 自然数|しぜんすう で は ない}{minus drei ist eine ganze Zahl, aber keine natürliche Zahl}
\lies{$\frac34 \in \Q$}{4|よん 分|ぶん の 3|さん は 有理数|ゆうりすう で ある}{drei Viertel ist eine rationale Zahl}
\lies{$\sqrt{2} \notin \Q$}{ルート 2|に は 無理数|むりすう で ある}{die Wurzel aus zwei ist irrational}
\lies{$0{,}\overline{3} = \frac13$ \newline \footnotesize jap. $0.\dot{3} = \frac13$}{循環|じゅんかん 小数|しょうすう 0.333|れい_てん_さん_さん_さん \ldots\ は 、 3|さん 分|ぶん の 1|いち に 等しい|ひとしい}{null Komma Periode drei ist gleich ein Drittel}
\end{lesen}

%%% Mengen
\begin{lesen}
\lies{$A = \{1;\,2;\,3;\,4\}$ \newline $B = \{3;\,4;\,5\}$}{$A$ は 1|いち 、 2|に 、 3|さん 、 4|よん を 要素|ようそ と する 集合|しゅうごう 、 $B$ は 3|さん 、 4|よん 、 5|ご を 要素|ようそ と する 集合|しゅうごう}{$A$ ist die Menge mit den Elementen eins, zwei, drei und vier; $B$ ist die Menge mit den Elementen drei, vier und fünf}
\lies{$A \cap B = \{3;\,4\}$}{$A$ と $B$ の 共通|きょうつう 部分|ぶぶん は 、 3|さん と 4|よん を 要素|ようそ と する 集合|しゅうごう}{die Schnittmenge von $A$ und $B$ ist die Menge mit den Elementen drei und vier}
\lies{$A \cup B = \{1;\,2;\,3;\,4;\,5\}$}{$A$ と $B$ の 和集合|わしゅうごう は 、 1|いち から 5|ご まで の 数|かず を 要素|ようそ と する}{die Vereinigungsmenge von $A$ und $B$ enthält die Zahlen von eins bis fünf}
\end{lesen}

%%% Gleichungen und Ungleichungen lösen
\begin{lesen}
\lies{$x^2 - 5x + 6 = 0$ \newline $(x - 2)(x - 3) = 0$}{左辺|さへん を 因数|いんすう 分解|ぶんかい する 。}{Zerlege die linke Seite in Faktoren.}
\lies{$x_1 = 2,\ x_2 = 3$}{積|せき が 0|ゼロ ならば 、 因数|いんすう の どれ か が 0|ゼロ で ある 。 解|かい は 2|に と 3|さん 。}{Ein Produkt ist null, wenn einer der Faktoren null ist. Die Lösungen sind zwei und drei.}
\end{lesen}

%%% Funktionen
\begin{lesen}
\lies{$f(x) = 2x^2 - 3$ \newline $f(2) = 5$}{エフ 2|に イコール 5|ご ： 点|てん $P(2,\,5)$ は $f$ の グラフ 上|じょう に ある}{$f$ von zwei ist fünf: Der Punkt $P(2 \mid 5)$ liegt auf dem Graphen von $f$}
\lies{$f(x) = \frac12 x + 1$}{傾き|かたむき 2|に 分|ぶん の 1|いち 、 $y$ 切片|せっぺん 1|いち の 一次|いちじ 関数|かんすう}{eine lineare Funktion mit der Steigung ein Halb und dem $y$-Achsenabschnitt eins}
\lies{$f(x) = 0$}{$f$ の 零点|れいてん を 求めなさい|もとめなさい 。 $f(x)$ が 0|ゼロ に なる $x$ の 値|あたい は いくつ か 。}{Bestimme die Nullstellen von $f$: Für welche $x$ ist $f$ von $x$ gleich null?}
\lies{$f(x) = g(x)$}{2つ|ふたつ の グラフ の 交点|こうてん を 求めなさい|もとめなさい 。}{Bestimme die Schnittpunkte der beiden Graphen.}
\end{lesen}

%%% Winkelfunktionen
\begin{lesen}
\lies{$\cos 60\degree = \dfrac12$}{コサイン 60|ろくじゅう 度|ど は 2|に 分|ぶん の 1|いち}{der Kosinus von sechzig Grad ist ein Halb}
\lies{$\sin x = 0{,}5$ \newline $x_1 = 30\degree,\ x_2 = 150\degree$}{0|れい 度|ど から 360|さんびゃくろくじゅう 度|ど まで の 範囲|はんい で 、 この 方程式|ほうていしき の 解|かい は 30|さんじゅう 度|ど と 150|ひゃくごじゅう 度|ど の 2つ|ふたつ で ある 。}{Für Winkel zwischen null und dreihundertsechzig Grad hat die Gleichung zwei Lösungen: dreißig Grad und hundertfünfzig Grad.}
\lies{$\sin(x + 360\degree) = \sin x$}{サイン 関数|かんすう は 、 周期|しゅうき 360|さんびゃくろくじゅう 度|ど の 周期|しゅうき 関数|かんすう で ある 。}{Die Sinusfunktion ist periodisch mit der Periode dreihundertsechzig Grad.}
\lies{$-1 \le \sin x \le 1$}{サイン 関数|かんすう の 値|あたい は 、 マイナス 1|いち 以上|いじょう 1|いち 以下|いか で ある 。}{Die Werte der Sinusfunktion liegen zwischen minus eins und eins.}
\end{lesen}

%%% Intervalle und Grenzwerte
\begin{lesen}
\lies{$x \in [2;\,5]$}{$x$ は 2|に 以上|いじょう 5|ご 以下|いか の 閉区間|へいくかん に 属する|ぞくする}{$x$ ist Element des abgeschlossenen Intervalls von zwei bis fünf}
\lies{$\lim\limits_{x \to \infty} \dfrac{1}{x} = 0$}{$x$ を 限りなく|かぎりなく 大きく|おおきく する と 、 $x$ 分|ぶん の 1|いち の 極限値|きょくげんち は 0|ゼロ}{der Grenzwert von eins durch $x$ für $x$ gegen unendlich ist null}
\lies{$\limn \dfrac{2n + 1}{n} = 2$}{$n$ を 限りなく|かぎりなく 大きく|おおきく する と 、 $n$ 分|ぶん の 2|に $n$ プラス 1|いち の 極限値|きょくげんち は 2|に}{der Grenzwert von zwei $n$ plus eins, geteilt durch $n$, für $n$ gegen unendlich ist zwei}
\lies{steigend: $2 \le x < \infty$}{関数|かんすう は 、 $x$ が 2|に 以上|いじょう の 範囲|はんい で 増加|ぞうか する}{die Funktion ist steigend für alle $x$, die größer oder gleich zwei sind}
\end{lesen}
